\section{How speech-to-speech systems work}
\label{sec:how}

There are three families of speech-to-speech conversational systems, and one important question for
each: \emph{does it exist for Amharic today, with open weights?}

\begin{figure}[H]
\centering
\begin{tikzpicture}[
  font=\sffamily\footnotesize, node distance=4mm,
  blk/.style={draw=ink2, rounded corners=2pt, minimum height=7mm, minimum width=15mm, align=center, fill=white},
  arr/.style={-{Stealth[length=2mm]}, ink2, semithick},
  lab/.style={font=\sffamily\small\bfseries, anchor=west, text=ink}]
% cascade
\node[lab] at (-0.2,0.9) {(a) Cascaded (pipeline)};
\node[blk] (m1) at (0.6,0) {mic};
\node[blk, right=of m1, fill=s1!12] (v1) {VAD /\\turn detect};
\node[blk, right=of v1, fill=s2!12] (a1) {ASR};
\node[blk, right=of a1, fill=s3!12] (l1) {text LLM};
\node[blk, right=of l1, fill=s4!15] (t1) {TTS};
\node[blk, right=of t1] (p1) {speaker};
\foreach \x/\y in {m1/v1,v1/a1,a1/l1,l1/t1,t1/p1} \draw[arr] (\x) -- (\y);
\node[anchor=west, text=good, font=\sffamily\scriptsize] at (p1.east) {\ \ Amharic: yes (all parts exist)};
% half cascade
\node[lab] at (-0.2,-1.3) {(b) Half-cascade (audio-in LLM)};
\node[blk] (m2) at (0.6,-2.2) {mic};
\node[blk, right=of m2, fill=s1!12] (v2) {VAD};
\node[blk, right=of v2, minimum width=37mm, fill=s3!12] (l2) {audio encoder + LLM\\{\scriptsize (e.g.\ Gemma 4 E2B/E4B)}};
\node[blk, right=of l2, fill=s4!15] (t2) {TTS};
\node[blk, right=of t2] (p2) {speaker};
\foreach \x/\y in {m2/v2,v2/l2,l2/t2,t2/p2} \draw[arr] (\x) -- (\y);
\node[anchor=west, text=warnc, font=\sffamily\scriptsize] at (p2.east) {\ \ Amharic: untested};
% end-to-end
\node[lab] at (-0.2,-3.5) {(c) End-to-end speech LLM (full duplex)};
\node[blk] (m3) at (0.6,-4.4) {mic};
\node[blk, right=of m3, minimum width=76mm, fill=s5!15] (e3) {speech tokens in $\rightarrow$ one model $\rightarrow$ speech tokens out\\{\scriptsize Moshi, GLM-4-Voice, Qwen3-Omni, LLaMA-Omni \dots}};
\node[blk, right=of e3] (p3) {speaker};
\draw[arr] (m3) -- (e3); \draw[arr] (e3) -- (p3);
\node[anchor=west, text=bad, font=\sffamily\scriptsize] at (p3.east) {\ \ Amharic: no};
\end{tikzpicture}
\caption{The three system families. Only the cascade can be built for Amharic with open weights
today. The half-cascade depends on whether an audio-input LLM understands Amharic speech, which is
still to be tested (experiment E-HC-1).}
\label{fig:families}
\end{figure}

\subsection{(a) Cascaded pipelines}
A voice-activity detector (VAD) finds speech; a turn detector decides that the user has finished;
automatic speech recognition (ASR) turns audio into text; a text LLM writes the reply; text-to-speech
(TTS) turns it back into audio. Every component can be swapped independently, and every stage can
be measured separately, which is exactly what the assignment requires. The cost is that latency
\emph{adds up}, unless the stages overlap.

\paragraph{Sequential vs.\ streaming cascades.}
The naive cascade waits for each stage to finish before starting the next. A \emph{streaming}
cascade passes partial results forward as soon as they exist:
\begin{itemize}
  \item \textbf{Streaming ASR} emits partial transcripts while the user is still talking, so at the
  endpoint only the last few hundred milliseconds of audio remain to decode. CTC encoders do this
  by processing overlapping chunks. Whisper can be forced into it with the \emph{LocalAgreement}
  policy~\cite{machacek2023}, at the cost of roughly 3\,s lag on GPU. Qwen3-ASR has a native
  2-second-chunk streaming mode~\cite{qwen3asr}.
  \item \textbf{Streaming LLM} output is standard: tokens arrive one by one. The important quantities
  are time to first token (TTFT, dominated by \emph{prefill} of the prompt) and time to the first
  \emph{speakable} chunk, meaning the first clause boundary.
  \item \textbf{Chunked or streaming TTS} synthesises the first clause while the LLM writes the next.
  Non-autoregressive models (VITS, diffusion) must wait for a whole chunk of text. Codec-LM TTS
  (e.g.\ Orpheus with SNAC) can emit audio frames incrementally.
\end{itemize}
Figure~\ref{fig:timeline} shows why this matters: in the streaming version, \ttfa\ contains only the
\emph{head} of each stage, not its full duration.

\begin{figure}[H]
\centering
\begin{tikzpicture}[font=\sffamily\scriptsize, x=0.82cm, y=0.62cm]
\def\row#1#2#3#4#5{\fill[#4] (#1,#2) rectangle (#3,#2+0.72); \node[anchor=center, text=ink] at ({(#1+#3)/2},#2+0.36) {#5};}
% axis
\draw[ink2, -{Stealth[length=1.6mm]}] (0,-0.4) -- (15.6,-0.4) node[anchor=west] {time};
% --- sequential
\node[anchor=east, font=\sffamily\small\bfseries] at (-0.1,4.7) {Sequential};
\row{0}{4.3}{3}{rule}{user speaks}
\row{3}{4.3}{3.8}{s1!55}{wait}
\row{3.8}{4.3}{5.3}{s2!60}{ASR}
\row{5.3}{4.3}{9.3}{s3!60}{LLM (whole reply)}
\row{9.3}{4.3}{12.3}{s4!70}{TTS (whole reply)}
\row{12.3}{4.3}{15.3}{ink2!25}{playback}
\draw[decorate, decoration={brace, amplitude=4pt, mirror}, ink] (3,4.15) -- (12.3,4.15) node[midway, below=4pt] {\ttfa\ (sequential)};
% --- streaming
\node[anchor=east, font=\sffamily\small\bfseries] at (-0.1,1.7) {Streaming};
\row{0}{1.3}{3}{rule}{user speaks}
\fill[s2!25] (0.6,2.15) rectangle (3,2.45); \node[anchor=east, text=ink2] at (0.5,2.3) {streaming ASR};
\row{3}{1.3}{3.8}{s1!55}{wait}
\row{3.8}{1.3}{4.2}{s2!60}{}
\row{4.2}{1.3}{5.6}{s3!60}{LLM}
\row{5.6}{1.3}{6.8}{s4!70}{TTS}
\row{6.8}{1.3}{9.8}{ink2!25}{playback, chunk 1}
\node[anchor=south, text=ink2, font=\sffamily\tiny, align=center] at (4.9,2.05) {first\\clause};
\node[anchor=south, text=ink2, font=\sffamily\tiny, align=center] at (6.2,2.05) {first\\chunk};
\fill[s3!30] (5.6,0.55) rectangle (9.0,0.95); \node[text=ink2] at (7.3,0.75) {LLM continues};
\fill[s4!35] (9.0,0.55) rectangle (9.8,0.95);
\draw[decorate, decoration={brace, amplitude=4pt, mirror}, ink] (3,1.15) -- (6.8,1.15) node[midway, below=4pt] {\ttfa};
\end{tikzpicture}
\caption{Same components, two orchestrations (schematic, not to scale). Streaming moves ASR work
before the endpoint and overlaps LLM and TTS generation with playback. \ttfa\ then depends on the
endpoint wait, the tail of ASR, LLM time to first clause and TTS time for that clause.}
\label{fig:timeline}
\end{figure}

\subsection{(b) Half-cascades: audio-in LLMs}
Some recent LLMs accept audio directly and can answer spoken questions without a separate ASR step.
Among open models, Gemma~4 E2B/E4B include a $\sim$300M-parameter audio encoder, support ASR and
speech translation for audio up to 30\,s, and were pre-trained on 140+ languages~\cite{gemma4card}.
The model card does not say whether its audio understanding covers Amharic. If it does, the
ASR stage disappears, and so does the separate transcript decode that is very slow for Amharic
(Section~\ref{sec:m-tok}). This is a high-value, low-cost test (E-HC-1).

\subsection{(c) End-to-end speech LLMs}
Full-duplex models such as Moshi~\cite{moshi} (about 200\,ms practical latency),
LLaMA-Omni~\cite{llamaomni}, GLM-4-Voice~\cite{glm4voice} and Qwen3-Omni~\cite{qwen3omni} model
speech tokens directly and can respond while the user is still talking. They represent the lowest
latency possible in principle, but none of them produces Amharic speech. SeamlessM4T~v2 accepts
Amharic speech but outputs Amharic only as \emph{text}, not speech~\cite{seamless}. Training such a
model for Amharic is well beyond a two-week budget. \textbf{Conclusion: the realistic architecture is
a streaming cascade, possibly a half-cascade, and the work is making every head of every stage
short.}

\subsection{Turn-taking: the stage everyone underestimates}
\label{sec:how-turn}
A VAD such as Silero~\cite{silero} tells speech from silence in under a millisecond per 32\,ms
frame, but it cannot tell a \emph{thinking pause} from the \emph{end of a turn}. Pure-silence
endpointers therefore wait a fixed \emph{hangover} (often 300--800\,ms~\cite{webrtcventures2026}),
and that wait is added directly to \ttfa. There are two ways to shorten it:
\begin{itemize}
  \item \textbf{Semantic end-of-turn models}, which judge from prosody or words whether the speaker
  is done. Smart Turn v3 uses a Whisper-tiny encoder with an 8M-parameter classifier on the last 8\,s
  of audio, and covers 14 languages (not Amharic)~\cite{smartturn}. Text-based detectors run a small
  LM on the partial transcript.
  \item \textbf{Speculative response}: start ASR finalisation and LLM prefill at the \emph{first}
  short pause, and throw the work away if the user continues. This hides the endpoint wait whenever
  the guess is right.
\end{itemize}

\subsection{Where production systems spend their time}
Industry latency budgets for (English) voice agents in 2026 put turn detection at 150--400\,ms,
streaming ASR finalisation at 100--300\,ms, LLM TTFT at 250--500\,ms and TTS first audio at
100--200\,ms, for a total of about 0.6--1.2\,s~\cite{webrtcventures2026,deepgram2026}. The two
consistently largest items are \textbf{turn-taking} and \textbf{LLM time to first token}. Prompt
caching of a fixed system prompt typically removes 200--400\,ms of TTFT~\cite{webrtcventures2026}.
Section~\ref{sec:measure} tests how these proportions change for Amharic on a laptop CPU.
